
%%%%%%%%%%%%%%%%%%%%%%%%%%%%%%%%%%%%%%%%%
% Cada capítulo empieza con \chapter
%%%%%%%%%%%%%%%%%%%%%%%%%%%%%%%%%%%%%%%%%
\chapter{Introducción}

%%%%%%%%%%%%%%%%%%%%%%%%%%%%%%%%%%%%%%%%%
% Luego introduce el texto qeu desees
%%%%%%%%%%%%%%%%%%%%%%%%%%%%%%%%%%%%%%%%%
\textbf{¡ Enhorabuena !}, ya sólo te queda realizar tu TFG para ser un Ingeniero, ¡ un último esfuerzo !

\textit{Modifica éste capítulo para redactar la introducción a tu trabajo.}

Te hemos preparado esta plantilla para hacerte más fácil el proceso de escribir el ``tomo''. No te preocupes en ningún momento del estilo visual, sólo preocúpate de que el contenido sea bueno y conciso.

Te dejamos algunos ejemplos de cómo usar \LaTeX a lo largo de tu proyecto.

%%%%%%%%%%%%%%%%%%%%%%%%%%%%%%%%%%%%%%%%%
% Tu capítulo puedes separarlo en secciones, subsecciones, subsubsecciones ...
%%%%%%%%%%%%%%%%%%%%%%%%%%%%%%%%%%%%%%%%%
\section{Objetivos}

Explica cuáles son los objetivos de tu propuesta. A lo largo del documento deberás explicar cómo los has resuelto y al final del documento, en el apartado  \nameref{apdo:conclusiones} (\ref{apdo:conclusiones})

\begin{itemize}
    \item Con el entorno `itemize' podrás poner un listado de \textit{bullets}
    \item Y cada \textbf{item} saldrá en una nueva línea
    \item También puedes ocupar varias líneas haciendo un salto con doble \textit{backslash} \\
        estoy en otra línea del mismo \textbf{bullet}.
    \item Y puedes agregar más subniveles si lo ves interesante, pero no abuses, haz que la lectura sea agradable.
    \begin{itemize}
        \item Soy un subnivel
    \end{itemize}
\end{itemize}

También puedes hacer listas numeradas:

\begin{enumerate}
    \item Soy el primero
    \item Soy el segundo
\end{enumerate}
